\subsubsection*{The AI Retailer (Single-Player Game)}

The bot is a \textbf{hybrid}: a deterministic economic solver makes every
\textit{decision}, while a \textbf{llama3} model served via Ollama (remote
hosts from the session config, with the local instance as failover) writes
the natural-language chat, and a dedicated \textit{offer-reader} model
(config key \texttt{llm\_reader}) extracts offers out of the human's
free-text messages.

\paragraph*{1. The setting}
A contract is a wholesale price $w \in [1,5]$\,€ (two decimals) and a
quantity $q \in \{1,\dots,100\}$ of widgets. Demand is $D \sim U[0,100]$
and expected sales for an order of $q$ are
$$ ES(q) \;\equiv\; E[\min(q,D)] \;=\; \frac{q(200-q)}{200}. $$
Expected profits (the Supplier carries the inventory risk):
\begin{itemize}
    \item Supplier: $\pi_S = w \cdot ES(q) - PC \cdot q$
    \item Retailer: $\pi_R = (RP - w) \cdot ES(q)$
\end{itemize}

Crucially, the bot's entire strategy runs on its \textbf{acting retail
price} --- the value it disclosed to the human, or the no-disclosure
fallback of 4€ --- never the true draw
(\texttt{bot\_acting\_market\_price}; the ``no full information'' design).

\paragraph*{2. The benchmark: Nash bargaining solution (\texttt{optimal.py})}
With $RP$ and $PC$, the symmetric even-split offer is:
$$ w^* = \frac{RP(RP + 3PC)}{2(RP + PC)}, \qquad
   q^* = 100\cdot\frac{RP - PC}{RP}
   \text{ (rounded to the total-profit-maximizing integer)} $$

The bot's \textbf{target profit} $\Pi^*$ is the supplier profit at
$(w^*, q^*)$, floored to two decimals. For $RP=4, PC=1$: $w^*=2.80$,
$q^*=75$, $\Pi^* = 56.25$.

\paragraph*{3. Offer evaluation (\texttt{offer.py})}
Every incoming offer (chat or interface) is classified:
\begin{itemize}
    \item \textbf{ACCEPT} iff the bot's own expected profit meets the
    target: $(RP - w)\cdot ES(q) \ge \Pi^*$.
    \item Otherwise, two \textbf{feasibility} tests diagnose what's wrong:
    \begin{itemize}
        \item \textbf{Price feasible} iff even at $q=100$ the bot could
        still hit the target: $(RP - w)\cdot ES(100) \ge \Pi^* \iff
        w \le RP - \Pi^*/50$. At $RP=4$: $w \le 2.875$, so the
        \textbf{price cap} is 2.87€.
        \item \textbf{Quantity feasible} iff some allowed price still hits
        the target at that $q$: $RP - \Pi^*/ES(q) \ge PC$. At $RP=4$ this
        requires $ES(q) \ge 18.75$, i.e. $q \ge 21$ --- the
        \textbf{minimum workable quantity} is 21 units.
    \end{itemize}
    \item Both fail $\rightarrow$ \texttt{NOT\_PROFITABLE\_ON\_BOTH};
    only price $\rightarrow$ \texttt{\_ON\_PRICE}; only quantity
    $\rightarrow$ \texttt{\_ON\_QUANTITY}.
    \item Out-of-range terms $\rightarrow$ \texttt{INVALID\_OFFER}; no
    readable terms $\rightarrow$ \texttt{NOT\_OFFER}.
\end{itemize}

\paragraph*{4. Mathematical answers to partial offers}
A chat message naming only one term becomes a partial offer, and the bot
\textbf{completes} it at its target profit:
\begin{itemize}
    \item \textbf{Quantity-only (``quantity 70'')} $\rightarrow$
    \texttt{OFFER\_QUANTITY} if $q \ge 21$. The bot keeps $q$ and solves
    $(RP - w)\cdot ES(q) = \Pi^*$ for the \textbf{highest} price it can
    grant, floored to cents:
    $$ w = \left\lfloor \left(RP - \frac{\Pi^*}{ES(q)}\right)\cdot 100
       \right\rfloor / 100 $$
    Examples at $RP=4$: $q=70 \Rightarrow ES=45.5 \Rightarrow w = 2.76$;
    $q=40 \Rightarrow 2.24$; $q=30 \Rightarrow 1.79$. Note this is
    \textbf{generous to the supplier}: any lower price also works for the
    bot, but it anchors on the boundary price.

    \item \textbf{Price-only (``how about 2.50?'')} $\rightarrow$
    \texttt{OFFER\_PRICE} if $w \le 2.87$. The bot keeps $w$ and picks the
    quantity that (1) satisfies its target and (2) \textbf{maximizes the
    supplier's profit} among candidates (total profit as tie-break). The
    feasible band comes from the quadratic
    $(RP-w)\cdot\frac{q(200-q)}{200} \ge \Pi^*$, whose roots are:
    $$ q_{\pm} = 100 \pm \sqrt{10000 - \frac{200\,\Pi^*}{RP - w}} $$
    Candidates are integers around $q_-$, $q_+$ (clamped to $[0,100]$) and
    around the retailer's vertex $q=100$; among the feasible ones it
    returns $\arg\max_q\, w\cdot ES(q) - PC\cdot q$. Examples:
    $w=2.50 \Rightarrow q=52$; $w=2.87 \Rightarrow q=94$.

    \item \textbf{Infeasible partials} ($q<21$, or $w>2.87$) $\rightarrow$
    treated as \texttt{NOT\_PROFITABLE\_ON\_BOTH} $\rightarrow$ the bot
    \textbf{anchors its counter on the Nash offer itself} (2.80€, 75).
    Same for unreadable messages (\texttt{NOT\_OFFER}). Invalid
    (out-of-range) offers get a reminder of the allowed ranges with no
    counter anchor. Complete-but-unprofitable offers use the same two
    solvers: the failing term is replaced while the workable term the
    human proposed is kept.
\end{itemize}

\paragraph*{5. The pipeline around the math}
Per chat message:
\begin{enumerate}
    \item \textbf{Guard classification} (\texttt{bot\_guard.py}, llama3),
    \textbf{context-aware}: the classifier sees the last three exchanged
    messages plus the bot's last line, so short contextual answers (a bare
    ``25'' after the bot asked for a quantity) are read as partial offers
    rather than noise, and their terms are backfilled if the offer reader
    misses them. Off-topic requests and non-widget goods $\rightarrow$
    fixed refusal line; verbal acceptances (``ok, deal'') $\rightarrow$
    scripted pointer to the CONFIRM button.
    \item \textbf{Offer reading}: spacy fast path, then the offer-reader
    Ollama model extracting [price, quantity] (range outputs like
    ``60--85'' resolve to the Nash terms when they contain them).
    \item \textbf{Decision} --- accept or reject, and the anchor terms of
    any counter --- by the math above, never by the LLM.
    \item \textbf{Reply generation} (\texttt{bot\_strategy.py}): llama3
    writes up to 3 candidate replies with the solver's anchor terms
    embedded in the prompt; each candidate is scrubbed, linted (no closure
    claims, no wrong-direction price talk, no code), and its offer is
    \textbf{read back out} and re-evaluated --- sent only if the read-back
    offer meets the bot's profit target, otherwise best-of-batch among
    lint-clean drafts, otherwise the solver-scripted fallback line. The
    binding counter posted to the interface is the offer read back out of
    the sent message. Accepted chat offers are re-posted as a binding
    interface offer for the human to CONFIRM; accepted interface offers
    close the deal directly.
\end{enumerate}

One asymmetry worth knowing when reading transcripts: the solver's anchor
terms always sit exactly \textbf{on} the bot's profit-target boundary ---
the supplier-favoring choice within that boundary is deliberate (it
maximizes the human's profit subject to the bot's target). The generated
reply may state terms that deviate from the anchor, but the read-back
validation only lets it deviate \textbf{above} the boundary (profit
$\ge \Pi^*$), never below.
